%% fig1_pipeline.tex -- Figure 1. The whole idea in one picture.
%%
%% FINAL FIGURE DESCRIPTION (for the illustrator, and for the record):
%% A left-to-right pipeline across the full text width, with a side panel.
%%   Left, stacked: the two trigger inputs, an operational alert and the demand
%%   telemetry stream. They meet in the trigger box, which is annotated with the
%%   refractory window of 5 periods and the cap of at most 2 proposals.
%%   The trigger opens a proposal opportunity at the stopping time tau_j.
%%   The language model is consulted exactly there and returns one typed
%%   ShockSpec, drawn with a lock icon to mark that it is frozen and immutable.
%%   The compiler maps the spec to one of 72 registered ControlConfig points.
%%   Between the compiler and the controller sits the e-process gate, drawn as a
%%   valve rather than an arrow, because that is the point of the figure: the
%%   compiled configuration is computed but does not flow until the evidence
%%   path E_{j,t}, accumulated only from observations after tau_j, crosses
%%   1/alpha_j. A dashed bypass path carries the baseline configuration straight
%%   to the controller and is the live path while the gate is closed.
%%   The controller emits q_t. A feedback arrow returns the realized observation
%%   Y_{t+1} to both the telemetry stream and the evidence path, which is what
%%   makes the closed-loop argument of Lemma 1 necessary rather than decorative.
%%   Right panel: the eleven-rung ladder as a vertical list, with arms 8, 9 and
%%   10 braced together and labeled "one physical call set, three activation
%%   policies", and arm 10 marked as the contribution.
%%
%% This file is the real figure, drawn in TikZ with base libraries only. It is
%% included by \input, not \includegraphics, so no external build step exists.

\begin{figure*}[!t]
\centering
\footnotesize
\begin{tikzpicture}[
  node distance=3.6mm,
  box/.style={draw, rounded corners=1pt, align=center, inner sep=2.2pt,
              minimum height=7.5mm, font=\footnotesize},
  src/.style={box, fill=black!4, minimum width=15mm},
  proc/.style={box, fill=black!6, minimum width=17mm},
  spec/.style={box, fill=black!10, minimum width=19mm, very thick},
  gate/.style={box, fill=black!14, minimum width=21mm, very thick},
  outbox/.style={box, fill=black!4, minimum width=11mm},
  rung/.style={align=left, font=\scriptsize, inner sep=1pt},
  flow/.style={-{Latex[length=1.6mm]}, thick},
  bypass/.style={-{Latex[length=1.6mm]}, thick, dashed},
  fb/.style={-{Latex[length=1.6mm]}, semithick, densely dotted},
  lbl/.style={font=\scriptsize, align=center, inner sep=1pt},
]

%% ---- inputs ---------------------------------------------------------------
\node[src] (alert) {operational\\alert};
\node[src, below=5mm of alert] (tele) {demand\\telemetry};

%% ---- trigger --------------------------------------------------------------
\node[proc, right=6mm of $(alert)!0.5!(tele)$] (trig) {trigger\\$\tau_j$};
\node[lbl, below=1.2mm of trig] {refractory 5\\$\le 2$ proposals};

%% ---- model ---------------------------------------------------------------
\node[proc, right=7mm of trig] (llm) {language\\model};

%% ---- the frozen spec ----------------------------------------------------
\node[spec, right=7mm of llm] (spec) {\shockspec{}\\\emph{frozen at} $\tau_j$};

%% ---- compiler ------------------------------------------------------------
\node[proc, right=7mm of spec] (comp) {compiler\\72 points};

%% ---- the gate ------------------------------------------------------------
\node[gate, right=7mm of comp] (gate)
  {\eprocess{} gate\\$E_{j,t} \ge 1/\alpha_j$};
\node[lbl, below=1.2mm of gate] {evidence from\\$\tau_j{+}1$ onward only};

%% ---- controller and order ------------------------------------------------
\node[proc, right=7mm of gate] (ctrl) {capped\\base-stock};
\node[outbox, right=6mm of ctrl] (order) {order $q_t$};

%% ---- flow ----------------------------------------------------------------
\draw[flow] (alert.east) -- (trig.west);
\draw[flow] (tele.east)  -- (trig.west);
\draw[flow] (trig)  -- (llm);
\draw[flow] (llm)   -- (spec);
\draw[flow] (spec)  -- (comp);
\draw[flow] (comp)  -- (gate);
\draw[flow] (gate)  -- (ctrl);
\draw[flow] (ctrl)  -- (order);

%% ---- the baseline bypass, live while the gate is closed ------------------
\coordinate (bypassL) at ($(comp.north) + (0,5.5mm)$);
\coordinate (bypassR) at ($(ctrl.north) + (0,5.5mm)$);
\draw[bypass] (comp.north) |- (bypassL) -- (bypassR) -| (ctrl.north);
\node[lbl, above=0.4mm] at ($(bypassL)!0.5!(bypassR)$)
  {baseline configuration, live until activation};

%% ---- closed loop ---------------------------------------------------------
\coordinate (fbY) at ($(tele.south) + (0,-8mm)$);
\draw[fb] (order.south) |- (fbY) -| (tele.south);
\node[lbl, above=0.4mm] at ($(fbY) + (26mm,0)$) {$Y_{t+1}$, and it also feeds the evidence path};

%% The ladder side panel was removed once Table~\ref{tab:ladder} carried the same
%% eleven rungs; keeping both duplicated the list and made this figure a third of
%% a page taller. The brace annotation that mattered is kept inline below.
\node[lbl, below=1.2mm of order] {arms 8, 9, 10 share\\one physical call set};

\end{tikzpicture}
\caption{The handshake. A shared trigger decides \emph{when} the model is
consulted, a closed schema decides \emph{what} it may say, and an
anytime-valid \eprocess{} decides \emph{whether} the answer touches an order.
The gate is drawn as a valve rather than an arrow because that is the point: the
compiled configuration is computed immediately and flows only after evidence
gathered strictly after $\tau_j$ drives $E_{j,t}$ above $1/\alpha_j$. The dashed
path is the baseline configuration, which is live until then. The dotted return
is what makes fact (F1) of Section~\ref{sec:verify} load-bearing rather than
decorative, since the controller's own orders shape the observations the test
consumes.}
\label{fig:pipeline}
\end{figure*}
